\section{Cơ chế tư vấn theo tình huống}
\label{sec:co-che}
\subsection{Luồng xử lý câu hỏi}
Từ kiến trúc ở Mục~\ref{sec:kien-truc}, một lượt hỏi đáp bắt đầu bằng việc
nhận diện mục đích, ghi nhận dữ kiện và xác định nội dung thuộc phạm vi hỗ
trợ. Hệ thống lựa chọn thẻ và nguồn ứng viên, kiểm tra căn cứ rồi xác định
phần có thể trả lời và phần cần hỏi bổ sung. Nghiên cứu về hỏi làm rõ trong
hội thoại tìm kiếm thông tin~\cite{aliannejadi2019clarifying} cung cấp cơ sở
cho việc lựa chọn câu hỏi tiếp nối thay vì chỉ tạo một câu trả lời cuối.
Trong đề tài, câu hỏi bổ sung được ràng buộc bởi dữ kiện có ảnh hưởng đến
điều kiện áp dụng, đáp ứng F01.

\begin{figure}[H]
\centering\begin{tikzpicture}[node distance=6mm and 9mm]
\node[block,text width=58mm] (start) {Nhận câu hỏi và cập nhật dữ kiện};
\node[block,text width=58mm,below=of start] (choose) {Chọn thẻ và nguồn ứng viên\\theo quyền, đối tượng và thời điểm};
\node[decision,below=of choose,text width=39mm] (source) {Đủ căn cứ để xét\\phần đang hỏi?};
\node[block,text width=36mm,right=of source] (stop) {Dừng phần thiếu căn cứ\\Nêu giới hạn, đầu mối};
\node[decision,below=of source,text width=39mm] (facts) {Thiếu dữ kiện\\làm đổi kết luận?};
\node[block,text width=36mm,right=of facts] (ask) {Hỏi bổ sung\\hoặc nêu khả năng có điều kiện};
\node[block,text width=58mm,below=of facts] (answer) {Đối chiếu và trả lời\\Dữ kiện, kết luận, căn cứ, bước tiếp};
\draw[arrow] (start)--(choose);
\draw[arrow] (choose)--(source);
\draw[arrow] (source)--node[above,font=\footnotesize]{Không}(stop);
\draw[arrow] (source)--node[left,font=\footnotesize]{Có}(facts);
\draw[arrow] (facts)--node[above,font=\footnotesize]{Có}(ask);
\draw[arrow] (facts)--node[left,font=\footnotesize]{Không}(answer);
\draw[arrow] (ask.east)--++(5mm,0)|-(start.east);
\end{tikzpicture}

\caption{Luồng trả lời, hỏi bổ sung và dừng ở phần thiếu căn cứ}
\label{fig:luong}
\end{figure}

Hình~\ref{fig:luong} thể hiện luồng đối với một phần câu hỏi. Với yêu cầu gồm
nhiều nội dung, từng phần được xét theo căn cứ riêng; phản hồi có thể đồng
thời đưa ra hướng dẫn và hỏi thông tin còn thiếu. Khi người dùng không cung
cấp dữ kiện, hệ thống nêu giới hạn hoặc các khả năng có điều kiện, giữ nguyên
những giá trị chưa xác định.

\subsection{Quản lý ngữ cảnh và cập nhật thông tin}
Theo dõi trạng thái hội thoại giúp duy trì thông tin theo miền và trường
qua các lượt trao đổi~\cite{wu2019trade}. Để phục vụ F03, thiết kế của hệ
thống tổ chức trạng thái tư vấn tại lượt $t$ thành:
\begin{equation}
S_t=(F_t,I_t,E_t,K_t),
\label{eq:state}
\end{equation}
trong đó $F_t$ là dữ kiện hiện tại, $I_t$ là tập thẻ đang xét, $E_t$ là căn
cứ cùng trạng thái áp dụng và $K_t$ là các kết luận với quan hệ phụ thuộc.
Biểu diễn này xác định thông tin cần quản lý trong thiết kế; các lựa chọn
lưu trữ cụ thể được thực hiện khi xây dựng ứng dụng.

Mỗi dữ kiện lưu giá trị, mức rõ ràng, lượt cung cấp và quan hệ thay thế nếu
được sửa. Thông tin cũ còn trong lịch sử để giải thích diễn tiến hội thoại,
nhưng không được sử dụng như giá trị hiện tại. Một phát ngôn chuyển sang
việc mới cần được phân biệt với việc sửa thông tin của tình huống đang xét.

Khi người dùng bổ sung, sửa hoặc rút lại dữ kiện, hệ thống chọn lại thẻ và
nguồn theo tình huống mới. Những hướng dẫn dựa trên dữ kiện hoặc căn cứ đã
thay đổi được xét lại; phần vẫn phù hợp được giữ và phần được thay thế cần
chỉ rõ lý do. Nếu người dùng nói ``đổi nguồn hỗ trợ'' mà chưa rõ thay hoàn
toàn hay kết hợp các nguồn, hệ thống hỏi làm rõ trước khi loại nguồn cũ.

\subsection{Đối chiếu điều kiện và lựa chọn căn cứ}
F02 yêu cầu các nguồn liên quan được áp dụng trên cùng một tình huống.
Chẳng hạn, điều kiện gắn với nhiệm vụ và quy định tài chính có thể cần đối
chiếu cùng hoạt động, đơn vị và nguồn kinh phí. Hình~\ref{fig:phu-thuoc} mô
tả quan hệ giữa dữ kiện, điều kiện, đoạn nguồn và kết luận. Việc lưu quan
hệ này giúp giải thích căn cứ và xác định phần cần xét lại khi thông tin đổi.

\begin{figure}[H]
\centering\begin{tikzpicture}[node distance=9mm and 18mm]
\node[block,text width=48mm] (facts) {Dữ kiện người dùng\\Hoạt động, nguồn, thời điểm};
\node[block,text width=48mm,right=of facts] (evidence) {Đoạn nguồn\\Điều kiện và phạm vi áp dụng};
\node[block,text width=55mm,below=12mm of $(facts.south)!.5!(evidence.south)$] (rule) {Điều kiện đã xác nhận\\Các dữ kiện quyết định};
\node[block,text width=55mm,below=of rule] (claim) {Kết luận từng phần\\Trạng thái, căn cứ và giới hạn};
\draw[arrow] (facts.south)--(rule.north west);
\draw[arrow] (evidence.south)--(rule.north east);
\draw[arrow] (rule)--(claim);
\end{tikzpicture}

\caption{Quan hệ dữ kiện, điều kiện, đoạn nguồn và kết luận}
\label{fig:phu-thuoc}
\end{figure}

Nguồn được lựa chọn theo đối tượng, đơn vị, thời điểm và điều khoản áp dụng.
Ngày ban hành, ngày hiệu lực, ngày sự việc và ngày hỏi cần được phân biệt.
Một nhiệm vụ được phê duyệt trước quy định mới có thể cần xét điều khoản
chuyển tiếp. Tương tự, giấy tờ tồn đọng không tự quyết định dùng quy định
cũ hay mới; mốc áp dụng phải được xác định từ nguồn liên quan.

Nếu người dùng chỉ nhớ tháng hoặc năm, dữ kiện đó có thể đủ khi không làm
phát sinh nhiều khả năng áp dụng khác nhau. Khi khoảng khai báo đi qua mốc
thay đổi quy định, hệ thống cần hỏi chính xác hơn hoặc nêu rõ giới hạn.
Hai đoạn nguồn khác nhau về đối tượng hoặc thời điểm cũng cần được phân
biệt với mâu thuẫn thực sự trong cùng một trường hợp.

\subsection{Tạo phản hồi và xử lý giới hạn}
Phản hồi tư vấn gồm dữ kiện đang dùng, kết luận theo từng phần, căn cứ hỗ
trợ, thông tin còn thiếu và bước tiếp theo. Độ dài phụ thuộc câu hỏi; câu
tra cứu đơn giản có thể được trả lời ngắn. Khi người dùng sửa thông tin,
phản hồi cần chỉ rõ hướng dẫn nào thay đổi và vì sao.

ALCE xem xét chất lượng câu trả lời có trích dẫn và mức hỗ trợ của căn cứ
được dẫn~\cite{gao2023alce}. Thiết kế vận dụng nguyên tắc này cho F05: mỗi
kết luận cần gắn với đoạn nguồn hỗ trợ đúng nội dung đang trả lời. Liên
kết đến tài liệu cần kèm điều, khoản hoặc đoạn đã sử dụng để người hỏi
kiểm tra được hướng dẫn.

Ba mức kết luận của F04 là \emph{đáp ứng theo khai báo}, \emph{không đáp ứng
theo khai báo} và \emph{chưa đủ căn cứ}. Chúng áp dụng cho từng điều kiện,
không xác nhận dữ kiện khai báo hoặc thay quyết định của đơn vị có thẩm
quyền. Nghiên cứu câu hỏi không có đáp án trong ngữ cảnh như
SQuAD~\cite{rajpurkar2018squad} cung cấp cơ sở cho việc xác định khi nào
nguồn không hỗ trợ một câu trả lời.

Thiếu dữ kiện có thể được giải quyết bằng hỏi bổ sung. Khi thiếu nguồn,
phạm vi chưa rõ hoặc còn mâu thuẫn chưa có cách giải quyết được xác nhận,
hệ thống dừng phần liên quan, giữ phần vẫn đủ căn cứ và hướng dẫn đầu mối
cần trao đổi. Với yêu cầu vượt phạm vi tư vấn, phản hồi giới hạn ở việc
giải thích và định hướng; bản tóm tắt để người dùng trao đổi với chuyên
viên không được tự chuyển thành hồ sơ hay gửi đi thay người hỏi.

\subsection{Minh họa hội thoại}
Xét tình huống hội thảo đã dùng ở Mục~\ref{sec:boi-canh}: người dùng dự kiến
sử dụng kinh phí nhiệm vụ, sau đó đổi sang hoàn toàn kinh phí đơn vị.
Bảng~\ref{tab:hoi-thoai} nối các cơ chế trên thành một quá trình hỏi đáp.
Ví dụ mô tả hành vi mong đợi của thiết kế; điều kiện và định mức cụ thể phải
được đối chiếu với nguồn đã xác nhận khi xây dựng nội dung ứng dụng.

\begin{table}[H]
\centering\small
\begin{tabularx}{\textwidth}{|P{16mm}|Y|Y|}
\hline
\textbf{Lượt} & \textbf{Thông tin người hỏi} & \textbf{Xử lý cần thiết} \\\hline
1 & Muốn dự hội thảo bằng kinh phí nhiệm vụ; chưa nêu thời điểm. & Ghi nhận hoạt động và nguồn; hỏi thời điểm hoặc dữ kiện khác chỉ khi ảnh hưởng. \\ \hline
2 & Bổ sung địa điểm, ngày hoạt động và quan hệ với nhiệm vụ. & Xét nguồn theo dữ kiện; trả phần có căn cứ, nêu phần chưa xác minh. \\ \hline
3 & Đổi sang hoàn toàn kinh phí đơn vị. & Chọn lại thẻ, nguồn; rút kết luận chỉ dành cho kinh phí nhiệm vụ. \\ \hline
4 & Quy định riêng của đơn vị chưa có trong kho. & Dừng phần cần nguồn riêng; giữ phần quy định chung còn đủ căn cứ và chỉ đầu mối. \\
\hline
\end{tabularx}
\caption{Minh họa cập nhật dữ kiện và căn cứ trong tình huống T6}
\label{tab:hoi-thoai}
\end{table}

Nguồn riêng chưa có ở lượt cuối là cách bố trí ví dụ để kiểm tra điểm dừng,
không phải nhận định HUST thiếu văn bản. Tình huống cho thấy việc đổi nguồn
kinh phí có thể dẫn tới chọn lại căn cứ và xét lại hướng dẫn; dữ kiện mới
không tự thay thế phần quy định đang thiếu trong kho.
